\documentclass{article} \usepackage[utf8]{inputenc} \usepackage[spanish]{babel} \usepackage{amsmath, amssymb, amsthm} \usepackage{bm} \usepackage{float} \usepackage{longtable} \usepackage[hidelinks]{hyperref}

\newtheorem{definition}{Definición} \newtheorem{theorem}{Teorema} \newtheorem{remark}{Observación} \newcounter{pseudocode}

\newcommand{\xvec}{\mathbf{x}} \newcommand{\thetavec}{\boldsymbol{\theta}} \newcommand{\Fvec}{\mathbf{F}} \newcommand{\R}{\mathbb{R}} \newcommand{\norm}[1]{\left\lVert#1\right\rVert} \newcommand{\loss}{\mathcal{L}} \newcommand{\dataset}{\mathcal{D}} \newcommand{\jac}{\mathbf{J}}

\title{Introducción} \author{Noel Pérez Calvo} \date{\today}

\begin{document}

\maketitle

\section*{Introducción} \label{sec:introduccion} El estudio de sistemas de ecuaciones no lineales paramétricos es una herramienta fundamental en la modelación de fenómenos científicos y de ingeniería \cite{gershenfeld2012nature}.  Desde la determinación de puntos de equilibrio en modelos dinámicos \cite{strogatz2018nonlinear} hasta el análisis de circuitos electrónicos \cite{horowitz2015art} o la cinética de reacciones químicas \cite{laidler1987chemical}, estos sistemas permiten representar cómo interactúan las variables de un proceso en función de ciertos parámetros.  En todos estos casos, conocer los valores de las incógnitas que anulan el sistema —sus ceros o soluciones— resulta imprescindible para entender el comportamiento del fenómeno modelado y tomar decisiones a partir de él \cite{allgower1990numerical}.

Más allá de obtener soluciones numéricas para valores concretos de los parámetros, interesa saber cómo varía cada solución cuando los parámetros cambian dentro de un rango determinado.  Esa dependencia funcional permite realizar análisis de sensibilidad \cite{saltelli2008global}, predecir el comportamiento del sistema ante nuevas condiciones \cite{forrester2008engineering} y evitar el costo computacional de resolver el sistema desde cero para cada nueva tupla de parámetros \cite{eldred2009formulations}.  En este contexto, el propósito ya no es un valor numérico puntual, sino una expresión analítica que describa cada solución en función de los parámetros.

Desde el punto de vista numérico, la resolución de sistemas no lineales ha sido estudiada en profundidad \cite{burden2010numerical}.  Los métodos iterativos clásicos \cite{nocedal2006numerical} permiten encontrar una solución aislada cuando se dispone de una aproximación inicial adecuada, pero no garantizan la obtención de todas las soluciones posibles \cite{meleshko1991nonlinear}.  Para abordar esta limitación se han desarrollado estrategias capaces de enumerar el conjunto completo de raíces reales en un dominio acotado, entre ellas las técnicas de deflación \cite{ojika1983deflation} y los algoritmos de búsqueda exhaustiva con garantías de completitud \cite{bartnik2024mdcm}.  Estos enfoques son especialmente útiles cuando se necesita construir un conjunto de datos que incluya todas las soluciones del sistema para distintos valores de los parámetros.

En un plano distinto, la regresión simbólica \cite{koza1992genetic} persigue descubrir expresiones matemáticas que relacionan variables de entrada con variables de salida sin imponer una forma funcional predefinida.  Herramientas recientes como \texttt{PySR} \cite{cranmer2023pysr}, \texttt{SymbolFit} \cite{diniz2025symbolfit} o \texttt{leaf} \cite{leaf2024} han automatizado este proceso con creciente eficiencia.  Paralelamente, se han explorado estrategias para manejar conjuntos de datos donde coexisten puntos que responden a funciones matemáticas diferentes \cite{ukorigho2023competitive}.  Estas últimas resultan especialmente pertinentes cuando se enfrenta un sistema paramétrico cuyas soluciones no pueden ser descritas por una única función, sino que cada una responde a una expresión matemática diferente.

Un intento de conectar el mundo numérico con el simbólico es el presentado en \cite{merlet2024mixing}, donde se utilizan redes neuronales entrenadas sobre soluciones numéricas de sistemas paramétricos para predecir aproximaciones iniciales que luego se refinan con métodos iterativos.  Sin embargo, esta estrategia no proporciona expresiones analíticas explícitas: las redes funcionan como una caja negra que ofrece una aproximación numérica, pero no revela la forma funcional de la solución.

A pesar de estos avances, cuando un sistema no lineal paramétrico admite varias soluciones para una misma tupla de parámetros, ni la regresión simbólica convencional ni las técnicas competitivas existentes ofrecen una vía directa para recuperar todas las dependencias funcionales a partir de un mismo conjunto de datos.  A diferencia de un conjunto etiquetado donde se conoce de antemano a qué solución pertenece cada punto, aquí el dato sólo indica que un cierto vector de incógnitas satisface el sistema para unos parámetros dados, pero no incluye información que permita distinguir de cuál de las posibles soluciones se trata.  Como consecuencia, puntos originados por distintas funciones conviven en el mismo conjunto sin que se pueda separarlos a priori, y el desafío consiste en caracterizar cada una de esas funciones sin conocer cuántas son ni qué forma tienen.

Teniendo en cuenta lo anterior, la pregunta que guía esta investigación es: ¿es posible, mediante una estrategia de regresión simbólica, recuperar expresiones analíticas de los ceros de un sistema de ecuaciones no lineales paramétrico?

El objeto de estudio se sitúa en la intersección del análisis numérico y el aprendizaje automático simbólico.  De manera más específica, se aborda la regresión simbólica en el contexto de los sistemas de ecuaciones no lineales paramétricos.  El interés se centra en aplicarla a conjuntos de datos generados mediante resolución numérica, con el propósito de separar y descubrir las múltiples funciones solución que en ellos conviven.

El objetivo general de este trabajo es diseñar e implementar un algoritmo que recupere automáticamente las expresiones analíticas de todas las soluciones de un sistema de ecuaciones no lineales paramétricas.  Para alcanzarlo, el algoritmo se apoya en una estrategia de regresión simbólica.  A partir de este objetivo se desprenden los siguientes objetivos específicos:

\begin{itemize} \item Construir un procedimiento automático que genere un conjunto de entrenamiento formado por tuplas (parámetros, solución), muestreando el espacio de parámetros y resolviendo numéricamente el sistema para cada tupla, de modo que todas las soluciones queden registradas.

\item Desarrollar una estrategia de regresión simbólica que, a partir del conjunto de datos completo, descubra secuencialmente las funciones que representan cada una de las soluciones del sistema, separando de manera automática las distintas dependencias funcionales.

\item Validar experimentalmente la metodología sobre sistemas de ecuaciones no lineales con número conocido de soluciones, comparando las expresiones recuperadas con las soluciones exactas cuando sea posible y analizando la precisión de las funciones obtenidas. \end{itemize}

La solución propuesta consiste en un algoritmo que alterna etapas de entrenamiento de modelos de regresión simbólica con un paso de filtrado de puntos ya explicados.  En cada iteración se entrena un modelo basado en programación genética sobre los datos aún no cubiertos; de entre todas las ecuaciones candidatas que produce el modelo, se selecciona la que consigue ajustar una mayor fracción de puntos y, a continuación, esos puntos se retiran del conjunto de trabajo.  El proceso se repite hasta que no quedan puntos suficientes.

Los experimentos realizados sobre sistemas con múltiples soluciones muestran que el algoritmo es capaz de recuperar las expresiones analíticas que dieron origen a los datos, incluso cuando las distintas soluciones se entremezclan en el conjunto de entrenamiento.  En los casos estudiados, las funciones fueron descubiertas con una cobertura superior al 95\,\% de los puntos, y su forma coincidió con las expresiones exactas conocidas.

La estructura del documento es la siguiente.  En el Capítulo 1 se presentan los fundamentos de los sistemas de ecuaciones no lineales paramétricos, los métodos numéricos para el cálculo de sus soluciones y los principios de la regresión simbólica basada en programación genética.  El Capítulo 2 describe la metodología propuesta: la generación de la malla de parámetros y del conjunto de entrenamiento, la estrategia iterativa de descubrimiento de funciones y los criterios de selección y parada.  También se detalla la arquitectura del sistema implementado, incluyendo los módulos de definición de ecuaciones, generación de la malla, resolución numérica, preparación de datos y regresión simbólica.  El Capítulo 3 expone los experimentos realizados y el análisis de los resultados.  Finalmente, se presentan las conclusiones del trabajo. \newpage \bibliographystyle{plain} \bibliography{referencias} \end{document}
